\title{FlashAttention-3: Fast and Accurate Attention with Asynchrony and Low-precision}

\begin{abstract}
Attention mechanisms form the critical computational layer within the widely adopted Transformer architecture, yet they remain the primary performance constraint for scaling large language models and accommodating extended input sequences. \textsc{FlashAttention} introduced a strategy to accelerate attention computation on GPUs by reducing the overhead of memory operations. Nonetheless, prior iterations have not leveraged enhancements in the latest hardware generations; for example, \textsc{FlashAttention-2} utilizes just 35\% of the available computational capacity on the H100 GPU. To address this limitation, we propose three key strategies to accelerate attention processing on Hopper GPUs: (1) harnessing the asynchronous nature of both Tensor Cores and TMA to enable warp-specialization that overlaps computation and data transfer, (2) fusing block-wise matrix multiplication with the softmax step by interleaving their execution, and (3) utilizing hardware-accelerated FP8 low-precision via block quantization and incoherent processing. Our new approach, \textsc{FlashAttention-3}, delivers 1.5-2.0$\times$ speedups on H100 GPUs—achieving performance of up to 740 TFLOPs/s (75\% hardware utilization) with FP16, and nearly 1.2 PFLOPs/s with FP8. Furthermore, our evaluation demonstrates that FP8 \textsc{FlashAttention-3} reduces numerical error by a factor of 2.6 compared to standard FP8 attention baselines.
\end{abstract}

\section{Introduction}
\label{sec:intro}

Within the Transformer framework~\citep{vaswani2017attention}, the self-attention module serves as the main computational constraint, as evaluating attention scores between queries and keys incurs a quadratic cost relative to the input sequence length. Enhancing attention mechanisms to support extended context lengths holds the potential to enable a variety of new functionalities—such as comprehending and reasoning over multiple long-form documents~\citep{guo2021longt5,shaham2022scrolls,peng2023yarn} and managing entire software repositories~\citep{roziere2023code, li2023starcoder}—as well as expanding to different data types like high-definition images~\citep{chen2022scaling}, audio signals~\citep{gulati2020conformer}, and video streams~\citep{ho2022video}. Moreover, applications such as tracking extended user histories~\citep{sun2019bert4rec} or planning in tasks that require long-term reasoning~\citep{yao2022react} can greatly benefit from these advancements. Consequently, there has been robust research interest in accelerating attention for longer sequences, through methods involving approximation~\citep{katharopoulos2020transformers,choromanski2020rethinking, tay2020efficient}, advances in software implementations (\citep{rabe2021self,dao2022flashattention,kwon2023efficient}), and investigations into alternative neural architectures~\citep{peng2023rwkv,sun2023retentive,gu2023mamba}.

In this study, we extend the contributions of \citet{dao2022flashattention}, who pioneered exact-attention algorithms that explicitly incorporate knowledge of GPU execution paradigms and hardware attributes into their overarching framework. Dao et al., in \citep{dao2022flashattention}, introduced \textsc{FlashAttention}, which leverages an innovative tiling approach for parallelizing attention computations. By consolidating all attention-related operations into a singular GPU kernel, this method avoids redundant accesses to global memory, thereby minimizing intermediate reads and writes.

Building upon this, \citet{dao2023flashattention2} presented \textsc{FlashAttention-2}, which reimagines the algorithm to facilitate parallelization across the sequence length axis. This variant processes blocks of key and value matrices in the inner loop of the forward pass, leading to enhanced workload allocation and GPU utilization. Despite these advances, our findings indicate that \textsc{FlashAttention-2} still demonstrates suboptimal utilization on recent GPU architectures compared to highly tuned matrix-multiplication (GEMM) kernels—achieving only around 35\% utilization on Hopper H100 GPUs, in contrast to the 80–90\% attained by GEMM kernels. This performance disparity can be partially explained by differences in implementation, such as the absence of Hopper-tailored instruction sets and continued reliance on instructions designed for Ampere architectures when employing Tensor Cores. Recent efforts, including ThunkerKitten~\citep{spector2024thunder} and cuDNN 9~\citep{cudnn9}, demonstrate that integrating Hopper-specific instructions along with tile-based programming paradigms can accelerate attention computations while also streamlining implementation.

At a deeper level, the algorithm implemented by \textsc{FlashAttention-2} operates within a straightforward, synchronous framework, lacking intentional integration of asynchronous execution or reduced-precision computation in its architecture. Asynchrony arises due to the hardware’s focus on expediting key machine learning operations, with specialized components such as Tensor Cores conducting matrix multiplications or Tensor Memory Accelerator (TMA) units handling memory access, functioning independently from standard CUDA cores responsible for logic, integer, and floating-point tasks. The adoption of lower numerical precision, including FP8 in Hopper and FP4 in Blackwell—building upon earlier use of FP16 (introduced with Pascal in 2017) and BF16 (with Ampere in 2020)—demonstrates a successful strategy to achieve up to twofold or fourfold increases in throughput at no extra cost in energy or silicon real estate. We provide an overview of Hopper’s contributions along these lines in \cref{subsec:hardware}. The principal technical obstacle is reformulating \textsc{FlashAttention-2} so that it can take full advantage of these hardware optimizations: leveraging asynchrony entails orchestrating the concurrent execution of matrix multiplication and softmax computations, despite their dependency relationship, and utilizing lower-precision arithmetic demands strategies to control quantization error—an especially pressing issue for rare or extreme values present in LLMs~\citep{dettmers2208llm, sun2024massive}.

In pursuit of this goal, we introduce \textsc{FlashAttention-3}, integrating and advancing three novel strategies to achieve superior efficiency on modern GPU hardware:\footnote{Although our experiments are grounded in NVIDIA’s Hopper architecture, the proposed method is broadly applicable to any GPU platform supporting efficient asynchronous execution and low-precision arithmetic.}

\begin{enumerate}

\item \textbf{Producer-Consumer asynchrony:} We introduce a warp-specialized software pipelining strategy that leverages the asynchronous capabilities of both data transfer and Tensor Core computations. By allocating the roles of data producers and consumers to distinct warps, this scheme enhances the algorithm’s potential to mask the latencies associated with memory accesses and instruction dispatch.

\item \textbf{Hiding softmax under asynchronous block-wise GEMMs:} Our approach overlaps the execution of lower-throughput operations needed for softmax (such as floating-point fused multiply-add and exponential functions) with the asynchronous block-wise GEMM operations performed by WGMMA instructions. In re-engineering the \textsc{FlashAttention-2} algorithm, we restructure it to avoid particular serial dependencies between softmax and GEMMs. For instance, in the two-stage variant of the algorithm, one block of the score matrix is processed by softmax concurrently with WGMMA prefetching and computing the next block asynchronously.

\item \textbf{Hardware-accelerated low-precision GEMM:} We modify the forward pass to utilize FP8 Tensor Cores for GEMM operations, resulting in almost a twofold increase in achieved TFLOPs/s. This adaptation addresses the different memory layout requirements for WGMMA regarding the organization of FP32 accumulators and FP8 operand matrices. To address the reduction in numerical precision associated with FP8, we incorporate block quantization and incoherent processing strategies, thereby lessening potential accuracy degradation.

In order to empirically assess our approach, we conduct benchmarks of \textsc{FlashAttention-3} on the H100 SXM5 GPU across varying configurations, demonstrating that (1) the FP16 version delivers a 1.5-2.0$\times$ speedup over \textsc{FlashAttention-2} during the forward computation (peaking at 740 TFLOPs/s), and a 1.5-1.75$\times$ improvement in the backward computation, (2) FP8 mode approaches 1.2 PFLOPs/s, and (3) at large sequence lengths, FP16 displays superior performance while FP8 maintains competitive throughput\footnote{To elaborate, with head dimension 64, \textsc{FlashAttention-3} FP8 outperforms, whereas for head dimensions of 128 and 256, it matches performance in non-causal configurations but trails when causal masking is enabled.} compared to the latest attention operator available in NVIDIA’s cuDNN library. Further, we confirm that FP16 \textsc{FlashAttention-3} matches \textsc{FlashAttention-2} in terms of numerical error and exhibits greater accuracy than the standard attention mechanism, as intermediate computations (such as those in softmax scaling) remain in FP32. In addition, FP8 \textsc{FlashAttention-3} employing block quantization and incoherent processing achieves 2.6$\times$ higher accuracy than conventional attention utilizing per-tensor quantization, particularly in scenarios characterized by outlier features.

We have released \textsc{FlashAttention-3} as open source under a permissive license\footnote{\textsc{FlashAttention-3}
  is available at \texttt{https://github.com/Dao-AILab/flash-attention}}, and intend to incorporate it into PyTorch and Hugging Face ecosystems to maximize accessibility and utility for the research and developer communities.

\section{Background: Multi-Head Attention and GPU Characteristics}
\label{sec:background}

\subsection{Multi-Head Attention}
\label{subsec:multi_head_attn}

Consider $\mathbf{Q}, \mathbf{K}, \mathbf{V} \in \mathbb{R}^{N \times d}$, representing the query, key, and value sequences for a single attention head, with $N$ denoting the length of the sequence and $d$ indicating the dimension of the head. The resulting attention output $\mathbf{O}$ is given by:

\begin{equation*}
  \mathbf{S} = \alpha \mathbf{Q} \mathbf{K}^\top \in \mathbb{R}^{N \times N}, \quad \mathbf{P} = \mathrm{softmax}(\mathbf{S}) \in \mathbb{R}^{N \times N}, \quad \mathbf{O} = \mathbf{P}\mathbf{V} \in \mathbb{R}^{N \times d},
\end{equation*}

Here, the $\mathrm{softmax}$ operation is performed across the rows, and it is common to choose the scaling coefficient $\alpha = 1/\sqrt{d}$. To avoid issues of numerical instability that arise from the exponentiation step, practitioners often subtract $\mathrm{rowmax}(\mathbf{S})$ from $\mathbf{S}$ prior to applying the softmax. In the context of multi-head attention (MHA), each attention head maintains its distinct query, key, and value projection matrices. This computation can be executed in parallel for several heads and input batches, resulting in the aggregate output tensor.

Let $\phi$ denote a scalar loss function, and define $\mathbf{d}(-) = \partial \phi / \partial (-)$ to represent its gradient with respect to the specified argument. Given $\mathbf{dO} \in \mathbb{R}^{N \times d}$, the gradients $\mathbf{dQ}$, $\mathbf{dK}$, and $\mathbf{dV}$ are determined via the chain rule as shown below:

\begin{align*}
  \mathbf{dV} &= \mathbf{P}^\top \mathbf{dO} \in \mathbb{R}^{N \times d} \\
  \mathbf{dP} &= \mathbf{dO} \mathbf{V}^\top \in \mathbb{R}^{N \times N} \\
  \mathbf{dS} &= \mathrm{dsoftmax} (\mathbf{dP}) \in \mathbb{R}^{N \times N} \\
  \mathbf{dQ} &= \alpha \mathbf{dS} \mathbf{K} \in \mathbb{R}^{N \times d} \\
  \mathbf{dK} &= \alpha \mathbf{dS}^\top \mathbf{Q} \in \mathbb{R}^{N \times d},
\end{align*}

In this context, $\mathbf{d}s = (\mathrm{diag}(p) - p p^\top)\mathbf{d}p$ expresses the relationship where $p = \mathrm{softmax}(s)$, treating $p$ as a function of the vector $s$. We denote the application of this expression to each row individually as $\mathrm{dsoftmax}(\mathbf{dP})$. Moreover, for the backward computation in MHA, this process can be efficiently parallelized over the batch dimension and the number of heads.

\subsection{GPU hardware characteristics and execution model}
\label{subsec:hardware}

We discuss the components of the GPU execution model that pertain to \textsc{FlashAttention-3}, centering our attention on the NVIDIA Hopper architecture as a specific realization of this framework.

\paragraph{Memory hierarchy:} The GPU features a layered organization of memory types, where higher-capacity memory comes at the cost of reduced bandwidth, as detailed in \cref{tab:gpu-hierarchy}\footnote{\citet{luo2024benchmarking} notes that shared memory achieves a bandwidth of 128 bytes per clock cycle per SM; we compute aggregate bandwidth by multiplying by 132 SMs and the maximum boost frequency of 1830 MHz.}. The most prominent memory, global memory (GMEM) or HBM, consists of off-chip DRAM that is accessible to all streaming multiprocessors (SMs). Information stored in GMEM is automatically brought into an on-chip L2 cache when accessed. Additionally, every SM is equipped with a relatively small, on-chip shared memory (SMEM), a highly banked cache that is managed explicitly by the programmer. Finally, each SM incorporates its own register file for the fastest data access.

\paragraph{Thread hierarchy:} The organization of the GPU's programming paradigm centers on structured groups of threads as execution entities. Arranged from the most granular to the broadest level, the thread hierarchy consists of individual threads, warps (comprising 32 threads each), warpgroups (which include 4 consecutive warps), threadblocks (also referred to as cooperative thread arrays or CTAs), threadblock clusters (specific to Hopper architectures), and, finally, grids.

The two hierarchies exhibit a strong interconnection. Threads belonging to a particular CTA are scheduled together on a single SM, and similarly, CTAs within the same cluster are scheduled concurrently on the same GPC. All threads within a CTA have direct access to SMEM, while each thread individually is allocated up to 256 private registers (RMEM).

\paragraph{Asynchrony and warp-specialization:}

GPUs are designed as throughput-oriented processors, leveraging high levels of concurrency and asynchronous operations to mask both memory and execution delays. In the Hopper architecture, asynchronous data transfers between GMEM and SMEM are managed by a specialized hardware component known as the Tensor Memory Accelerator (TMA) \cite[\S7.29]{cuda}. Moreover, distinguishing itself from earlier generations like Ampere, Hopper’s Tensor Core—accessible through the warpgroup-wide WGMMA instruction \cite[\S9.7.14]{ptx}—supports asynchronous execution and is capable of directly reading inputs from shared memory.

The provision of hardware-level asynchrony enables warp-specialized kernel designs, in which the warps within a CTA are assigned exclusively to either producer or consumer tasks—executing solely data transfers or computational operations, respectively. Such specialization generally enhances the compiler’s capacity to optimize instruction scheduling \citep{warp-specialization-2011}. Furthermore, the Hopper architecture introduces support for dynamic register allocation among warpgroups through the \verb|setmaxnreg| instruction \cite[\S9.7.17.1]{ptx}. This feature permits warps engaged in matrix multiply-accumulate (MMA) operations to be allocated a greater portion of RMEM, whereas warps responsible for TMA (which requires only one thread) can utilize fewer registers.

\paragraph{Low-precision number formats:}
\label{sec:low-precision-gpu}
Recent advancements in GPU architectures include dedicated hardware components designed to boost the efficiency of low-precision arithmetic operations. As an illustration, Hopper’s WGMMA instruction is optimized for FP8 Tensor Cores, offering twice the performance per SM relative to FP16 or BF16 formats.

Effectively utilizing FP8 WGMMA requires careful attention to the operand layout requirements. Consider a GEMM operation performing $A \times B^{\top}$, where $A$ has dimensions $M \times K$ and $B$ has dimensions $N \times K$. The $A$ or $B$ operand is referred to as \emph{mn-major} if it exhibits memory contiguity along the outer $M$ or $N$ dimension, whereas it is termed \emph{k-major} if it is stored contiguously along the inner $K$-dimension. In the case of FP16 WGMMA, both mn-major and k-major arrangements are permissible for SMEM-resident operands. In contrast, FP8 WGMMA restricts operand layout to the k-major format exclusively. Furthermore, challenges arise in applications like attention mechanisms, where chaining multiple GEMMs within a single kernel is desirable—here, the incompatibility between FP32 accumulator layout and FP8 operand organization complicates the correct sequencing of dependent FP8 WGMMA invocations.

Within the realm of attention mechanisms, these constraints on layout necessitate specific adjustments to how an FP8 algorithm is constructed, as detailed in \cref{sec:algofp8}.

\subsection{Standard Attention and Flash Attention} In line with the definition from \citet{dao2022flashattention}, we refer to \textbf{standard attention} as a GPU-based attention computation that writes out the intermediate matrices $\mathbf{S}$ and $\mathbf{P}$ to HBM. The core contribution of \textsc{FlashAttention} was the adoption of a localized softmax reduction strategy, thereby eliminating costly intermediate memory operations and merging the attention computation into a unified kernel. The local softmax mechanism is implemented in lines \ref{code-ws:softmax_start}-\ref{code-ws:softmax_end} of the consumer mainloop within \cref{alg:flash3_wgmma_ws_only}, including the rescaling steps applied to blocks of $\mathbf{O}$. A straightforward justification that this approach accurately produces $\mathbf{O}$ is provided in \cite[\S 2.3.1]{dao2023flashattention2}.

\section{FlashAttention-3: Algorithm}
\label{sec:algo}

In this section, we present the details of the \textsc{FlashAttention-3} algorithm. For clarity, we concentrate on the forward computation, deferring the explanation of the backward computation to~\cref{sec:algo_ws_bwd}. We begin by outlining how to incorporate warp-specialization together with a circular SMEM buffer into the foundational framework of \textsc{FlashAttention-2}. Next, we detail the approach for leveraging WGMMA asynchrony to establish a two-stage pipeline that overlaps GEMM and softmax operations. Lastly, we discuss the adjustments required for supporting FP8, addressing both layout compatibility and accuracy improvements using block quantization and incoherent computation techniques.

\subsection{Producer-Consumer asynchrony through warp-specialization and pingpong scheduling}
\label{sec:algo_ws}

\paragraph{Warp-specialization}
Similar to the approach used in \textsc{FlashAttention-2}, the forward computation in \textsc{FlashAttention-3} exhibits substantial parallelism across dimensions such as batch size, head count, and the length of the query sequence. Therefore, it is sufficient to analyze the algorithm at the CTA level, where it processes a submatrix $\mathbf{Q}_i$ from the query matrix and computes the corresponding output submatrix $\mathbf{O}_i$. For clarity, we initially present the warp-specialization mechanism employing a circular SMEM buffer, omitting, for now, any overlapping between GEMM and softmax computations. Define $d$ as the head dimension, $N$ as the total sequence length, and select a block size $B_r$ for the query. This partitions $\mathbf{Q}$ into $T_r = \lceil \frac{N}{B_r} \rceil$ blocks denoted as $\mathbf{Q}_1, .., \mathbf{Q}_{T_r}$.

\begin{algorithm}[H]
    \caption{\small\label{alg:flash3_wgmma_ws_only}\textsc{FlashAttention-3} forward pass \textbf{excluding} intra-consumer overlapping -- CTA perspective}
    \begin{algorithmic}[1]
\REQUIRE Matrices $\mathbf{Q}_i \in \mathbb{R}^{B_r \times d}$ and $\mathbf{K}, \mathbf{V} \in \mathbb{R}^{N \times d}$ stored in HBM, a key block size $B_c$ with $T_c = \lceil \frac{N}{B_c} \rceil$.
\STATE Create a pipeline construct to coordinate barrier synchronization, employing a circular shared memory buffer with $s$ stages.
\IF {producer warpgroup active}
\STATE Free a fixed quantity of registers as specified.
\STATE Initiate data transfer of $\mathbf{Q}_i$ from HBM into shared memory.
\STATE Once data load is finalized, signal to consumers that $\mathbf{Q}_i$ is available.
\FOR{$0 \le j < T_c$}
    \STATE Pause until the $(j\,\%\,s)$-th stage of the buffer is processed by consumers.
    \STATE Begin loading $\mathbf{K}_j, \mathbf{V}_j$ from HBM into shared memory at the $(j\,\%\,s)$-th buffer stage.
    \STATE After transfer finishes, notify consumer groups that $\mathbf{K}_j, \mathbf{V}_j$ are ready.
\ENDFOR
\ELSE \STATE Allocate the predefined register count based on the consumer warp configuration.
\STATE Initialize on-chip: $\mathbf{O}_i = (0) \in \mathbb{R}^{B_r \times d}$, and scalars $\ell_i, m_i = (0), (-\infty) \in \mathbb{R}^{B_r}$.
\STATE Wait until $\mathbf{Q}_i$ has been loaded to shared memory.
\FOR{$0 \le j < T_c$}
\STATE Await completion of $\mathbf{K}_j$ loading in shared memory.
\STATE Calculate $\mathbf{S}_i^{(j)} = \mathbf{Q}_i \mathbf{K}_j^T$ (SS-GEMM); commit computation and synchronize. \label{alg:ws_only_gemm1}
\STATE Record previous value $m_i^{\mathrm{old}} = m_i$ and update $m_i = \mathrm{max}(m_i^{\mathrm{old}}, \mathrm{rowmax}(\mathbf{S}_i^{(j)}))$. \label{code-ws:softmax_start}
\STATE Determine $\widetilde{\mathbf{P}}_i^{(j)} = \mathrm{exp}(\mathbf{S}_i^{(j)} - m_i)$ and update $\ell_i = \mathrm{exp}(m_i^{\mathrm{old}} - m_i) \ell_i + \mathrm{rowsum}(\widetilde{\mathbf{P}}_i^{(j)})$. \label{code-ws:softmax_end}
\STATE Await readiness of $\mathbf{V}_j$ in shared memory.
\STATE Update $\mathbf{O}_i = \mathrm{diag}(\mathrm{exp}(m_i^{\mathrm{old}} - m_i))^{-1} \mathbf{O}_i + \widetilde{\mathbf{P}}_i^{(j)} \mathbf{V}_j$ (RS-GEMM); commit changes and synchronize. \label{alg:ws_only_gemm2}
\STATE Release ownership of the $(j\,\%\,s)$-th buffer stage back to producer.
\ENDFOR
\STATE Finalize $\mathbf{O}_i = \mathrm{diag}(\ell_i)^{-1} \mathbf{O}_i$ and compute $L_i = m_i + \log(\ell_i)$.
\STATE Store $\mathbf{O}_i$ and $L_i$ in HBM, marking them as the $i$th segments of $\mathbf{O}$ and $L$.
\ENDIF
\end{algorithmic}
\end{algorithm}

In our deployment of \cref{alg:flash3_wgmma_ws_only} on Hopper, (de)allocations are managed via \verb|setmaxnreg|, while TMA is utilized to load $\mathbf{Q}_i$ as well as the set $\{ \mathbf{K}_j, \mathbf{V}_j \}_{0 \leq j < T_c}$. GEMMs within the consumer mainloop are performed using WGMMA, employing either the SS or RS prefix to indicate whether the first operand is obtained from shared memory or the register file, respectively. To understand the operational sequence of \cref{alg:flash3_wgmma_ws_only}, observe that TMA load operations can be dispatched asynchronously and do not block based on the completion status of preceding loads. Furthermore, within the producer mainloop, the initial $s$ iterations do not trigger any waits, as the buffer is still being populated.

\paragraph{Pingpong scheduling} The asynchrony between WGMMA and TMA, in conjunction with warp specialization, allows for the possibility of overlapping the softmax operation of one warpgroup with the GEMM computation of a different warpgroup. This approach is motivated by the observation that, on state-of-the-art hardware accelerators, non-matrix multiplication instructions execute at significantly lower throughput compared to matmul operations. For instance, on the H100 SXM5 GPU, FP16 matrix multiplications can reach 989 TFLOPS, while specialized operations like the exponential function\footnote{According to the CUDA programming guide, each streaming multiprocessor (SM) supports 16 special function operations per clock. Multiplying this by 132 SMs and a 1830 MHz clock yields a theoretical throughput of 3.9 TFLOPS for special functions.} (required in softmax) are limited to 3.9 TFLOPS. Considering an FP16 attention forward pass with a head dimension of 128, the computational load measured in FLOPS for matmuls is 512 times greater than for exponential evaluations, while the throughput for exponentials is 256 times smaller; consequently, exponentials may consume approximately 50\% of the total execution time relative to matmuls. The disparity becomes more pronounced with FP8: as matmul throughput doubles, the capability for exponential operations remains unchanged.

The computation of the exponential function utilizes a dedicated hardware unit, known as the multi-function unit. For maximum efficiency, it is preferable to schedule exponential operations concurrently with the matmul operations executed by the Tensor Cores. To achieve this concurrency, we implement synchronization barriers (\texttt{bar.sync} instructions), which enforce an execution order such that the GEMMs—specifically, GEMM1 (\(\mathbf{P} \mathbf{V}\) from the current iteration) and GEMM0 (\(\mathbf{Q} \mathbf{K}^\top\) from the subsequent iteration)—for warpgroup 1 are executed prior to those in warpgroup 2. This ordering ensures that the softmax computation for warpgroup 1 takes place simultaneously with the GEMMs of warpgroup 2. Subsequently, the process alternates, assigning the softmax calculation to warpgroup 2 while warpgroup 1 handles its GEMMs, resulting in what is referred to as “pingpong” scheduling. This mechanism is depicted in~\cref{fig:pingpong_scheduling}. Although the pingpong execution does not perfectly mirror the schematic in practice, we typically observe substantial performance improvements (for example, FP16 forward throughput increases from 570 TFLOPS to between 620 and 640 TFLOPS when using a head dimension of 128 and a sequence length of 8192).

\paragraph{Attention variants} In handling multi-query attention~\citep{shazeer2019fast} and grouped query attention~\citep{ainslie2023gqa}, we adopt the strategy used in \textsc{FlashAttention-2}, modifying the tensor indexing so as to prevent redundant storage of $\mathbf{K}$ and $\mathbf{V}$ within HBM.

\subsection{Intra-warpgroup overlapping GEMMs and softmax}

It is possible to interleave certain instructions from the softmax operation with those from the GEMMs, even within a single warpgroup. Below, we outline a specific method to achieve this overlap.

Within the attention mechanism, certain operations in the primary loop exhibit sequential dependencies that hinder parallel execution during any single iteration. Specifically, the (local) softmax computation (see lines \ref{code-ws:softmax_start}–\ref{code-ws:softmax_end}) depends on the output $\mathbf{S}_i^{(j)}$ generated by the initial GEMM, while the subsequent GEMM requires as input the value $\widetilde{\mathbf{P}}_i^{(j)}$ produced by the softmax. Notably, the synchronization points at lines \ref{alg:ws_only_gemm1} and \ref{alg:ws_only_gemm2} in \cref{alg:flash3_wgmma_ws_only} enforce a serial ordering between the softmax and GEMM operations. To circumvent these dependencies, we introduce a pipeline across loop iterations, employing supplemental register-resident buffers. Building on this approach, we introduce a two-stage\footnote{It is important to recognize that while the overlapping scheme can be configured for up to $s$ stages—matching the number of stages in the circular SMEM buffer—it is not limited to nor must it always utilize all $s$ stages.} pipelined GEMM-softmax algorithm as follows:

\begin{algorithm}[H]
    \caption{\small\label{alg:flash3_wgmma}\textsc{FlashAttention-3} consumer warpgroup forward pass}
    \begin{algorithmic}[1]
      \REQUIRE  Input matrices $\mathbf{Q}_i \in \mathbb{R}^{B_r \times d}$ along with $\mathbf{K}, \mathbf{V} \in \mathbb{R}^{N \times d}$ residing in HBM, a key block size $B_c$, and the number of blocks $T_c = \lceil \frac{N}{B_c} \rceil$.
      \STATE Adjust allocation of registers depending on the consumer warpgroup count.
      \STATE Locally, set $\mathbf{O}_i = (0) \in \mathbb{R}^{B_r \times d}$ as well as initialize $\ell_i = 0$, $m_i = -\infty$ in $\mathbb{R}^{B_r}$.
      \STATE Ensure both $\mathbf{Q}_i$ and the initial block $\mathbf{K}_0$ are present in shared memory before proceeding.
      \STATE With WGMMA, perform the calculation $\mathbf{S}_{\mathrm{cur}} = \mathbf{Q}_i \mathbf{K}_0^T$. Execute, then wait for completion.
      \STATE Release the first segment of the $\mathbf{K}$ buffer.
      \STATE Derive $m_{i}$, $\tilde{\mathbf{P}}_{\mathrm{cur}}$, and $\ell_{i}$ from $\mathbf{S}_{\mathrm{cur}}$, and apply rescaling to $\mathbf{O}_i$.
      \FOR{$1 \le j < T_c - 1$}
        \STATE Await $\mathbf{K}_j$ availability in shared memory.
        \label{code:mainloop_start}
        \STATE Compute $\mathbf{S}_{\mathrm{next}} = \mathbf{Q}_i \mathbf{K}_{j}^T$ through WGMMA. Issue the command but do not wait for completion. \label{code:first_wgmma}
        \STATE Await $\mathbf{V}_{j-1}$ to be loaded into shared memory.
        \STATE Update $\mathbf{O}_{i} = \mathbf{O}_{i} + \tilde{\mathbf{P}}_{\mathrm{cur}} \mathbf{V}_{j-1}$ using WGMMA. Execute the computation but do not block on it. \label{code:second_wgmma}
        \STATE Wait for WGMMA operation corresponding to $\mathbf{Q}_i \mathbf{K}_{j}^T$ to complete.
        \STATE Using $\mathbf{S}_{\mathrm{next}}$, calculate $m_{i}$, $\tilde{\mathbf{P}}_{\mathrm{next}}$, and $\ell_{i}$. \label{code:softmax}
        \STATE After the WGMMA result for $\tilde{\mathbf{P}}_{\mathrm{cur}} \mathbf{V}_{j-1}$ is ready, rescale $\mathbf{O}_i$
        \STATE Free the $(j\,\%\,s)$-th and $(j-1\,\%\,s)$-th stages in the $\mathbf{K}$ and $\mathbf{V}$ buffers, respectively.
        \STATE Move contents of $\mathbf{S}_{\mathrm{next}}$ to $\mathbf{S}_{\mathrm{cur}}$.
        \label{code:mainloop_end}
      \ENDFOR \STATE Wait for the loading of $\mathbf{V}_{T_c - 1}$ into shared memory.
      \STATE Calculate
      $\mathbf{O}_{i} = \mathbf{O}_{i} + \tilde{\mathbf{P}}_{\mathrm{last}} \mathbf{V}_{T_c - 1}$ by means of WGMMA, and ensure completion by waiting.

\STATE Epilogue: Adjust $\mathbf{O}_{i}$ according to $m_{i}$. Determine $L_{i}$ using $m_{i}$ and $\ell_{i}$.
Store $\mathbf{O}_{i}$ and $L_{i}$ in HBM as the $i$-th segments of $\mathbf{O}$ and $L$, respectively.
\end{algorithmic}
\end{algorithm}

\cref{alg:flash3_wgmma} serves to substitute the consumer path found in \cref{alg:flash3_wgmma_ws_only}, thereby providing the full implementation of the \textsc{FlashAttention-3} algorithm when using FP16 precision. In broad terms, the term WGMMA is employed as a shorthand for asynchronous GEMM. Throughout the mainloop, spanning lines \ref{code:mainloop_start} to \ref{code:mainloop_end}, the second invocation of WGMMA in iteration $j$ (line \ref{code:second_wgmma}) is executed concurrently with the softmax computations associated with iteration $j+1$ (line \ref{code:softmax}).

Although the pipelined approach depicted earlier promises performance improvements in theory, several real-world considerations must be taken into account:

\paragraph{Compiler reordering} The idealized order of execution in the pseudocode does not always reflect actual instruction flow, as the NVCC compiler may reorder instructions to enhance optimization. Such rearrangement can disturb the planned sequence of WGMMA and accompanying non-WGMMA operations, potentially undermining the intended overlap and reducing performance benefits. Nevertheless, a review of the generated SASS code confirms that overlapped execution generally aligns with expectations (Section~\ref{sec:2-stage-sass}).

\paragraph{Register pressure} Achieving high performance requires careful management of register usage to prevent register spilling. The introduction of a 2-stage pipeline brings added demand for registers, as it necessitates storing not only intermediate computation outcomes but also pipeline state. Notably, each threadblock must retain an additional $\mathbf{S}_{\mathrm{next}}$ in registers, incurring a cost of $B_r \times B_c \times \text{sizeof}(\text{float})$ in extra register storage. This heightened register usage can be at odds with increasing threadblock size, which also intensifies register demands. Effective optimization therefore hinges on empirical tuning informed by profiling data to find an appropriate balance.

\paragraph{3-stage pipelining} Building upon the 2-stage structure, we explore a 3-stage pipeline that further interleaves execution by overlapping a second WGMMA operation with the softmax step. While this deeper pipeline may allow for superior Tensor Core utilization, it also imposes an even greater requirement for registers due to the extra computation stage. This further complicates decisions regarding the optimal interplay between pipeline depth and tile dimensions. Details regarding the 3-stage design and its performance assessment appear in ~\cref{sec:3-stage}.

\subsection{Low-precision with FP8}
\label{sec:algofp8}

\textbf{Efficiency: layout transformations.} Executing the forward computation of \textsc{FlashAttention-3} using FP8 precision introduces extra difficulties regarding layout alignment, which are not present when operating in FP16.

To begin with, the input tensors $\mathbf{Q}$, $\mathbf{K}$, and $\mathbf{V}$ are conventionally organized to be contiguous along the head dimension. However, for the second GEMM operation under the FP8 WGMMA k-major format, the constraint is that $\mathbf{V}$—specifically, its tiles as they are staged into SMEM—must instead be contiguous with respect to the sequence length dimension. As the TMA loading mechanism does not alter the dimension along which data is contiguous, we are left with two alternatives: (1) performing a transpose of $\mathbf{V}$ within GMEM as a preprocessing measure, or (2) executing a transpose of $\mathbf{V}$'s tiles in-kernel after their transfer to SMEM. Pursuing the first strategy, we may (1a) combine the transpose into the epilogue of a previous operation like the rotary embedding, or (1b) invoke a dedicated pre-processing transpose kernel\footnote{A well-optimized transpose kernel can approach the theoretical device bandwidth~\citep{colfax_cutlass_transpose_2024}.} to swap the head and sequence length strides. Nonetheless, option (1a) poses integration challenges for typical libraries, while option (1b) is excessively inefficient in scenarios that are limited by memory bandwidth, as is often the case in inference workloads.

Rather, for FP8 \textsc{FlashAttention-3}, we choose the second strategy. During the in-kernel transpose, we leverage the LDSM (\verb|ldmatrix|) and STSM (\verb|stmatrix|) instructions, which enable a warp of threads to cooperatively move data between SMEM and RMEM at a 128-byte granularity.\footnote{According to the PTX documentation, LDSM and STSM operations copy $8 \times 8$ matrices with 16-bit elements \cite[\S 9.7.13.4.15-16]{ptx}; nonetheless, by packing two 8-bit values into each 16-bit element, these instructions can be applied for FP8 data. However, as the transpose variants of LDSM/STSM do not unpack individual 8-bit elements, we must perform certain register shuffling between the LDSM and STSM stages to effectuate the desired tile-wise transpose, though the specifics are omitted here.} LDSM and STSM offer efficient register utilization, making it feasible to execute them in the producer warpgroup while also supporting layout transpositions during memory operations. Additionally, beginning with the second iteration, we can pipeline the transposition of the subsequent $\mathbf{V}$ tile to occur concurrently with the two WGMMAs that process the prior $\mathbf{V}$ and the current $\mathbf{K}$ tile.

Second, it is notable that the FP32 accumulator utilized by FP8 WGMMA does not share the same memory organization as operand A when stored in registers, contrary to the situation in FP16. Illustrative examples of these differing layouts are provided in \cref{fig:rmem_accum} and \cref{fig:rmem_operand}, which indicate the sequence in which register entries are assigned per thread. By employing byte permute instructions, it becomes possible to reformat the output of the initial WGMMA's accumulator so that it aligns both with the requirements of a subsequent WGMMA and with the format of the $\mathbf{V}$ tile produced by the in-kernel transpose operation. To be specific, and as shown in \cref{fig:rmem_accum}, the bytes are reordered as follows:
$$\{ \verb|d0 d1 d4 d5 d2 d3 d6 d7| \},$$
with this pattern repeating every 8 bytes in the register file. From the perspective of the logical structure of the $\mathbf{P}$ tile, this procedure effectively reorders the columns, such that, for example, columns $0189$ are mapped to the initial four columns. To ensure WGMMA generates the appropriate output tile, the in-kernel transpose can be configured to emit a corresponding permutation of the rows within the $\mathbf{V}$ tile.\footnote{Leveraging this flexibility in the in-kernel transpose negates the requirement for additional shuffle instructions to transfer register contents between threads, a necessity previously discussed in~\citep{colfax_fp8_flashattention_2024}.}

\textbf{Accuracy: block quantization and incoherent processing.} The FP8 (e4m3) format allocates only 3 bits to represent the mantissa and 4 bits for the exponent. As a consequence, this format introduces greater numerical errors relative to FP16 or BF16. In addition, large-scale models often display a prevalence of outlier values~\citep{dettmers2208llm,
  sun2024massive}—entries whose magnitudes significantly exceed the majority—posing challenges for effective quantization. Standard practice involves per-tensor scaling~\citep{micikevicius2022fp8}, in which each tensor (such as $\mathbf{Q}$, $\mathbf{K}$, and $\mathbf{V}$) is associated with a single scale factor. 
To improve the numerical robustness of attention computations in FP8, we adopt two approaches:

\begin{enumerate}

\item \textbf{Block quantization}: In this approach, a single scalar is retained for each block, partitioning the tensors $\mathbf{Q}$, $\mathbf{K}$, and $\mathbf{V}$ into separate blocks of either size $B_r \times d$ or $B_c \times d$, and independently quantizing each block. This quantization step can be efficiently combined with the pre-attention operation, such as the application of rotary embedding, incurring no additional latency because the latter is limited by memory bandwidth rather than computation. Since \textsc{FlashAttention-3} inherently processes data in blocks, it is straightforward to rescale each corresponding block in $\mathbf{S}$ to incorporate the quantization without incurring further computational overhead.

\item \textbf{Incoherent processing}: To mitigate the effect of outliers, before applying FP8 quantization, we left-multiply $\mathbf{Q}$ and $\mathbf{K}$ by a random orthogonal matrix $\mathbf{M}$. The orthogonality of $\mathbf{M}$ ensures $\mathbf{M} \mathbf{M}^\top = I$, so $(\mathbf{Q} \mathbf{M})(\mathbf{K} \mathbf{M})^\top$ is equivalent to $\mathbf{Q} \mathbf{K}^\top$, leaving the attention calculation unchanged. This transformation helps diffuse the outlier effects, since each component of $\mathbf{Q} \mathbf{M}$ or $\mathbf{K} \mathbf{M}$ represents a stochastic combination of the original entries of $\mathbf{Q}$ or $\mathbf{K}$, thus decreasing quantization errors. In line with \citet{chee2024quip} and~\citet{tseng2024quip}, we implement $\mathbf{M}$ as a product of random diagonal matrices populated with $\pm 1$ entries and a Hadamard matrix. This construction enables efficient multiplication in $O(d \log d)$ time as opposed to $O(d^2)$, and can also be seamlessly integrated with the rotary embedding operation at no additional computational expense.
\end{enumerate}

We demonstrate in \cref{sec:numerical_error} that employing both of these methods leads to a reduction in numerical errors by as much as a factor of 2.6.

\section{Empirical Validation}
\label{sec:exp}

Our implementation and assessment of \textsc{FlashAttention-3}'s performance and correctness rely on CUTLASS~\citep{Thakkar_CUTLASS_2023} primitives, including the WGMMA and TMA abstractions.

\begin{itemize}

\item \textbf{Attention Benchmarking.} We evaluate the execution time of \textsc{FlashAttention-3} over varying sequence lengths, making direct comparisons to the baseline PyTorch implementation, \textsc{FlashAttention-2}, the Triton version of \textsc{FlashAttention-2} (leveraging H100-specific instruction sets), as well as a specialized \textsc{FlashAttention-2} implementation from cuDNN designed for H100 GPUs. Our results show that \textsc{FlashAttention-3} achieves up to a 2.0$\times$ speedup over \textsc{FlashAttention-2}, and is 1.5$\times$ faster than its Triton counterpart. Furthermore, \textsc{FlashAttention-3} attains peak throughput of 740 TFLOPs/s, representing 75\% of the peak theoretical TFLOPs/s available on H100 GPUs.

\item \textbf{Ablation Analysis.} We demonstrate that incorporating algorithmic enhancements—specifically, warp-specialization and the integration of GEMM-softmax pipelining—results in measurable performance gains for \textsc{FlashAttention-3}.

\item \textbf{FP8 Attention Accuracy.} Through experimentation, we show that applying block quantization combined with incoherent processing diminishes the numerical error in FP8 \textsc{FlashAttention-3} by a factor of 2.6$\times$.
\end{itemize}

\subsection{Benchmarking Attention}
\label{subsec:benchmark_attn}

We evaluate the execution time of various attention mechanisms using an H100 80GB SXM5 GPU under multiple configurations (with and without causal masking, and with head dimensions set to 64 or 128), utilizing FP16 data types. The outcomes, summarized in~\cref{fig:benchmark_attn_fwd} and~\cref{fig:benchmark_attn_bwd}, indicate that \textsc{FlashAttention-3} achieves a speedup of approximately 1.5 to 2.0$\times$ over \textsc{FlashAttention-2} during the forward computation and about 1.5 to 1.75$\times$ during the backward computation. When contrasted with a conventional attention baseline, \textsc{FlashAttention-3} demonstrates up to 3 to 16$\times$ improvements in performance. Furthermore, for medium and large sequence lengths (1k tokens and above), \textsc{FlashAttention-3} outperforms even the proprietary cuDNN library, which has been specially tuned for H100 GPUs.

\paragraph{Benchmark settings:} The sequence length is adjusted across values of 512, 1k, up to 16k, while the batch size is selected such that the total token count remains fixed at 16k. The hidden dimension is consistently assigned a value of 2048. For the head dimension, values of 64, 128, or 256 are chosen, which correspond to 32, 16, or 8 attention heads, respectively. The computation of forward pass FLOPs is performed using:

\begin{equation*}
  4 \cdot \text{seqlen}^2 \cdot \text{head dimension} \cdot \text{number of heads}.
\end{equation*}

Applying causal masking reduces the computation by about half, as only roughly 50% of the elements are evaluated, so we divide the value by 2. For the backward pass, the FLOPs are estimated by scaling the forward pass FLOPs by a factor of 2.5. This factor arises because, while the forward computation involves 2 matrix multiplications, the backward phase—accounting for recomputation—requires 5 matrix multiplications in total.

Additionally, we evaluate the forward pass runtime using FP8 under comparable conditions. The outcomes for headdim 256 are displayed in ~\cref{fig:benchmark_attn_fp8}, with comprehensive results provided in \cref{sec:benchmark_attn_fp8_full}.

\subsection{Ablation Study: 2-Stage Pipelining Experiments}
\label{sec:2-stage-pipelining-experiments}

We evaluate the impact of both the two-stage WGMMA-softmax pipelining and the use of warp-specialization within the non-causal FP16 \textsc{FlashAttention-3} kernel, using the parameter configuration $\{\text{batch}, \text{seqlen}, \text{nheads}, \text{hdim}\} = \{ 4, 8448, 16, 128\}$. As shown in~\cref{table:ablation_pipelining}, these algorithmic enhancements—specifically, the incorporation of asynchrony through warp-specialization and the concurrent execution of GEMM with softmax—yield a marked increase in throughput, raising performance from 570 to 661 TFLOPs.

\subsection{Numerical Error Validation}
\label{sec:numerical_error}

Given the ongoing discussion surrounding the numerical accuracy of \textsc{FlashAttention}~\citep{golden2024flash}, we evaluate \textsc{FlashAttention-2}, \textsc{FlashAttention-3}, and a conventional attention implementation, using a high-precision FP64 reference as the baseline. To model the presence of outlier activations and features typical in LLMs~\citep{dettmers2208llm, sun2024massive}, we sample the elements of $\mathbf{Q}, \mathbf{K}, \mathbf{V}$ from the distribution described below:

\begin{equation*}
  \mathcal{N}(0, 1) + \mathcal{N}(0, 100) \cdot \mathrm{Bernoulli}(0.001).
\end{equation*}

Specifically, every element is sampled from a normal distribution centered at zero with a standard deviation of 1; however, for 0.1\% of the elements, we superimpose an independent Gaussian noise term with a standard deviation of 10. The resulting root mean squared error (RMSE) is reported in~\cref{table:numerical_error}. Under FP16 precision, both \textsc{FlashAttention-2} and \textsc{FlashAttention-3} produce RMSE values that are 1.7$\times$ lower than those from the conventional method, as intermediate calculations (specifically, the softmax operation) remain in FP32. The baseline FP8 attention leverages per-tensor scaling, with matrix multiplication accumulation in FP32 and the softmax intermediates retained in FP16. Owing to the use of block quantization and asynchronous computation, \textsc{FlashAttention-3} operating in FP8 achieves a 2.6$\times$ improvement in accuracy relative to this baseline.

\section{Dicussion, Limitations, Conclusion}
\label{sec:discussion}

Through the development of \textsc{FlashAttention-3}, we have shown that leveraging innovative programming strategies and harnessing advanced hardware capabilities like asynchrony and reduced-precision computation can significantly enhance both the performance and fidelity of attention mechanisms. Our approach accelerates attention operations by 1.5-2.0$\times$ over \textsc{FlashAttention-2}, and it also achieves a 2.6$\times$ reduction in FP8 numerical error relative to conventional per-tensor quantization methods. Nevertheless, several aspects remain for future exploration, such as refining the implementation for LLM inference workloads, adapting the FP8 kernel to include a persistent kernel architecture,\footnote{In our experiments, the FP16 variant of \textsc{FlashAttention-3} incorporates both persistent kernel and load balancing features, whereas these are absent in the FP8 variant. This partially accounts for the diminished performance of FP8 \textsc{FlashAttention-3} on tasks involving short sequence lengths and causal masking, compared to the FP8 cuDNN alternatives.} and systematically investigating the influence of low-precision attention during large-scale model training. Although the primary focus of our study is on Hopper GPUs, we anticipate that these methods can be extended to a broader set of hardware accelerators. Ultimately, we believe that improvements in attention—both in speed and precision—will foster the development of innovative applications for long-context scenarios.

\end{document}